% Figure: linear strain and linear stress distributions across a beam
% cross-section in bending. The strain varies linearly from compression
% above the neutral axis to tension below; the stress follows by Hooke's
% law. The cross-section is shown at left, the strain profile centre, the
% stress profile right, all sharing the neutral axis.
\begin{figure}[htbp]
\centering
\begin{tikzpicture}[scale=1.0]
% cross-section at left
\draw[engblack, very thick] (0,-1.4) rectangle (0.9,1.4);
\node[engblack, anchor=south, font=\scriptsize] at (0.45,1.45) {section};
\draw[enggray, dashed] (-0.3,0) -- (1.2,0)
  node[anchor=west, font=\scriptsize]{N.A.};
\draw[<->, enggray] (1.05,0) -- (1.05,1.4)
  node[midway, anchor=west, font=\scriptsize]{$y_{\max}$};
% strain profile (centre)
\begin{scope}[xshift=3.0cm]
  \draw[enggray, ->] (-1.3,0) -- (1.5,0)
    node[anchor=west, font=\scriptsize]{$\epsilon$};
  \draw[enggray] (0,-1.5) -- (0,1.5);
  \draw[engblue, very thick] (1.0,-1.4) -- (-1.0,1.4);
  \fill[engblue!12] (0,-1.4) -- (1.0,-1.4) -- (0,0) -- cycle;
  \fill[engblue!12] (0,1.4) -- (-1.0,1.4) -- (0,0) -- cycle;
  \node[engblue, anchor=south, font=\scriptsize] at (-0.7,1.0) {comp.};
  \node[engblue, anchor=north, font=\scriptsize] at (0.7,-1.0) {tens.};
  \node[engblack, anchor=south, font=\scriptsize] at (0,1.55)
    {$\epsilon=-y/\rho$};
\end{scope}
% stress profile (right)
\begin{scope}[xshift=6.4cm]
  \draw[enggray, ->] (-1.3,0) -- (1.5,0)
    node[anchor=west, font=\scriptsize]{$\sigma$};
  \draw[enggray] (0,-1.5) -- (0,1.5);
  \draw[engred, very thick] (1.0,-1.4) -- (-1.0,1.4);
  \fill[engred!12] (0,-1.4) -- (1.0,-1.4) -- (0,0) -- cycle;
  \fill[engred!12] (0,1.4) -- (-1.0,1.4) -- (0,0) -- cycle;
  \node[engred, anchor=north, font=\scriptsize] at (0.85,-1.0)
    {$\sigma_{\max}$};
  \node[engblack, anchor=south, font=\scriptsize] at (0,1.55)
    {$\sigma=My/I$};
\end{scope}
\end{tikzpicture}
\caption[Linear strain and stress across a bending section]{The
plane-sections assumption forces the bending strain to vary linearly
across the depth, $\epsilon=-y/\rho$, zero at the neutral axis and
largest in magnitude at the extreme fibres. Hooke's law turns the linear
strain into a linear stress $\sigma=E\epsilon=My/I$, compression on one
side of the neutral axis and tension on the other. The resultant of the
stress over the area is zero, which places the neutral axis at the
centroid; the moment of the stress about the neutral axis is the bending
moment $M$.}
\label{fig:vol03ch08:strain-stress-distribution}
\end{figure}
